\documentclass[a4paper]{article}
\usepackage{amsmath}
\usepackage{amsfonts}
\usepackage{amssymb}
\usepackage{graphicx}
\usepackage{float}
\begin{document}
\noindent \large Auteur: Siebe Haché \\
\noindent \large Studierichting: Informatica \\
\noindent \large Oefeningenreeks 7E nr. 3.2 \\
\medskip
\normalsize
Bepaal van volgende functies van drie veranderlijken de kritieke punten, en ga na of de functie in die punten een extremum bereikt.\\

De gegeven functie is $f(x,y,z) = 2x^{2} + 6y^{2} + 4z^{2} - 4xy + 4xz$.\\

$\dfrac{\partial f}{\partial x} = 4x - 4y + 4z$\\

$\dfrac{\partial f}{\partial y} = 12y - 4x$\\

$\dfrac{\partial f}{\partial z} = 8z + 4x$\\

Kritieke punten:\\

$\left\{ \begin{array}{l}
4x - 4y + 4z = 0 \\
-4x + 12y = 0 \\
4x + 8z = 0
\end{array} \right.$\\

Uit de tweede vergelijking: $-4x + 12y = 0 \implies x = 3y$.\\

Uit de derde vergelijking: $4x + 8z = 0 \implies z = -\dfrac{x}{2} = -\dfrac{3y}{2}$.\\

Substitutie in de eerste vergelijking:\\

$4(3y) - 4y + 4\left(-\dfrac{3y}{2}\right) = 12y - 4y - 6y = 2y = 0 \implies y = 0$\\

Dus $y = 0$, $x = 0$, $z = 0$.\\

Het enige kritieke punt is $(0, 0, 0)$.\\

Tweede-orde partiële afgeleiden:\\

$\dfrac{\partial^{2} f}{\partial x^{2}} = 4$\\

$\dfrac{\partial^{2} f}{\partial y^{2}} = 12$\\ 

$\dfrac{\partial^{2} f}{\partial z^{2}} = 8$\\

$\dfrac{\partial^{2} f}{\partial x \partial y} = -4$\\

$\dfrac{\partial^{2} f}{\partial x \partial z} = 4$\\ 

$\dfrac{\partial^{2} f}{\partial y \partial z} = 0$\\

De Hessiaanse matrix in $(0,0,0)$ is:\\

$H(0) = \left(\begin{array}{ccc}
4 & -4 & 4 \\ -4 & 12 & 0 \\ 4 & 0 & 8 \end{array} \right)$\\

$H - \lambda I = \left( \begin{array}{ccc} 4-\lambda & -4 & 4 \\ -4 & 12-\lambda & 0 \\ 4 & 0 & 8-\lambda \end{array} \right)$\\

$\det(H - \lambda I) = (4-\lambda)\begin{vmatrix} 12-\lambda & 0 \\ 0 & 8-\lambda \end{vmatrix}
-(-4)\begin{vmatrix} -4 & 0 \\ 4 & 8-\lambda \end{vmatrix}
+4\begin{vmatrix} -4 & 12-\lambda \\ 4 & 0 \end{vmatrix}$\\

$p(\lambda) = \lambda^{3} - 24\lambda^{2} + 144\lambda - 64$\\\

Stel $\lambda \leq 0$. Dan geldt term per term:\\

$\lambda^{3} \leq 0, -24\lambda^{2} \leq 0,  144\lambda \leq 0,  -64 < 0$\\

Alle termen zijn dus strikt negatief of nul, zodat $p(\lambda) < 0$ voor alle $\lambda \leq 0$.\\

Dit betekent dat $p$ geen wortels kan hebben in $(-\infty, 0]$, en bijgevolg zijn alle eigenwaarden van $H(0)$ strikt positief.\\

Dus de functie $f$ bereikt in het kritieke punt $(0, 0, 0)$ een strikt lokaal minimum.\\

\begin{thebibliography}{999}
\end{thebibliography}
\end{document}